% =============================================================================
% CVE-2026-60004 — Xpectra Security Research
% Assembled by build.sh: cat main.tex part2.tex part3.tex part4.tex > paper-full.tex
% Every number here comes from designs/FACTS.md or a cited evidence file.
% =============================================================================
% !TEX program = pdflatex
\documentclass[journal]{IEEEtran}
\usepackage[utf8]{inputenc}
\usepackage{graphicx}
\usepackage{booktabs}
\usepackage{array}
\usepackage{xcolor}
\usepackage{listings}
\usepackage{hyperref}
\usepackage{microtype}
\usepackage{subcaption}

\hypersetup{colorlinks=true,linkcolor=blue,citecolor=blue,urlcolor=blue}

\definecolor{codebg}{RGB}{245,247,250}
\definecolor{codekw}{RGB}{0,90,150}
\definecolor{codecm}{RGB}{120,120,120}
\definecolor{codestr}{RGB}{160,40,40}

\lstdefinestyle{gosrc}{
  language=Go,
  backgroundcolor=\color{codebg},
  basicstyle=\ttfamily\scriptsize,
  keywordstyle=\color{codekw}\bfseries,
  commentstyle=\color{codecm}\itshape,
  stringstyle=\color{codestr},
  frame=single,
  breaklines=true,
  breakatwhitespace=false,
  columns=fullflexible,
  keepspaces=true,
  showstringspaces=false,
  tabsize=2,
  xleftmargin=0.5em, xrightmargin=0.3em
}
\lstdefinestyle{diffsrc}{
  language={},
  backgroundcolor=\color{codebg},
  basicstyle=\ttfamily\scriptsize,
  commentstyle=\color{codecm}\itshape,
  frame=single,
  breaklines=true,
  breakatwhitespace=false,
  columns=fullflexible,
  keepspaces=true,
  showstringspaces=false,
  tabsize=2,
  xleftmargin=0.5em, xrightmargin=0.3em
}
\lstdefinestyle{transcript}{
  language={},
  backgroundcolor=\color{codebg},
  basicstyle=\ttfamily\scriptsize,
  frame=single,
  breaklines=true,
  breakatwhitespace=false,
  columns=fullflexible,
  keepspaces=true,
  showstringspaces=false,
  xleftmargin=0.5em, xrightmargin=0.3em
}

\newcommand{\figdir}{figures}

\title{Beyond the Published Exploit: Process-Level Proof, Live-Negative Fix
Verification, and the Detection Consequences of Gitea's \texttt{diffpatch} RCE
(CVE-2026-60004)}

% JOURNAL MODE: \IEEEauthorblockN/A are inline dummies in journal mode — stack manually.
\author{Xpectra Security Research\\[0.5ex]
\small Autonomous Red/Blue Operations}

\markboth{Beyond the Published Exploit: CVE-2026-60004}{Xpectra Security
Research \MakeLowercase\expandafter\the\year}

\maketitle

\begin{abstract}
Gitea's \texttt{diffpatch} API endpoint allowed arbitrary shell execution as
the service operating-system account (CVE-2026-60004, CWE-94, affected
\texttt{>=1.17}, \texttt{<1.27.1}). The vendor advisory, published 2026-07-28,
did not merely describe the flaw: it shipped a working Python exploit of
roughly 330 lines, including an exfiltration technique that needs no outbound
channel. \textbf{This paper makes no discovery claim.} It contributes what that
advisory left unproven: process-level proof of execution captured from an
independent reimplementation (parent process \texttt{git write-tree}, hook
process \texttt{/bin/sh\hspace{0pt}hooks/post-index-change\hspace{0pt}0\hspace{0pt}0},
working directory inside the bare temporary clone, \texttt{GIT\_DIR=.},
\texttt{uid=1000(git)}); empirical verification of the three preconditions the
vendor only asserts, including the finding that all three hold in the official
Docker image without setting any precondition-related option (the only
container parameters were standard hostname, port, and install-lock values
equivalent to a standard install's own, none part of the precondition set);
an A/B fix verification whose
negative arm is live rather than vacuous (22 of 22 cycles without execution,
with arrival of the payload at the sink path positively observed in 2 of 22
sampled cycles, and the second \texttt{diffpatch} response changing from HTTP 201
to HTTP 500); a measured
impact radius (\texttt{app.ini} of 1323 bytes with 3 secret-bearing lines read,
planted secrets read back, arbitrary write under \texttt{/data/gitea}
confirmed, and one negative result); and an analysis of the fix whose decisive
hunk changes one boolean argument and leaves the \texttt{git apply} flags
untouched --- so a byte-identical patch is inert after the fix, which makes
patch-content-based detection structurally weak and forces behavioural
detection. All experiments ran in an isolated, zero-egress Docker lab against a
fictional internal domain; no third-party system was touched. CISA added the
CVE to the KEV catalog on 2026-08-25 with a federal deadline of 2026-08-28.
\end{abstract}

\begin{IEEEkeywords}
Gitea, CVE-2026-60004, remote code execution, Git hooks, vulnerability
verification, detection engineering, reproducible experiments.
\end{IEEEkeywords}

% =============================================================================
\section{Introduction}
\label{sec:intro}

A source-hosting service that applies attacker-supplied patches has to decide
where those patches land before they are ever committed. Gitea 1.27.0 applied
them inside a \emph{bare} temporary clone, and in a bare repository the
checkout root \emph{is} \texttt{\$GIT\_DIR}. A patch that adds an executable
file at \texttt{hooks/post-index-change} therefore does not drop a stray file:
it installs a live Git hook that Git then executes while writing the index.
That is CVE-2026-60004 --- arbitrary command execution as the Gitea service
user, from one authenticated, write-capable API call, sent twice.

This CVE is unusual in a way that shapes every sentence that follows:
\textbf{the vendor published the working exploit inside the advisory}
(GHSA-rcr6-4jqh-j84m, published 2026-07-28T17:34:46Z). The advisory carries the
full mechanism, the vulnerable command line, a complete Python
proof-of-concept of roughly 330 lines, and an exfiltration technique that
stores command output in Git objects and retrieves it as a branch through
authenticated smart HTTP --- no outbound connection required. We make no
discovery claim anywhere in this paper, in our video material, or in any
companion artifact; Section~\ref{sec:record} states the division of credit
explicitly.

What an advisory cannot do for you is the rest of the security work:
\emph{showing that it actually happens, showing that the fix actually stops
it, and showing what a defender can actually see}. Working in an isolated
zero-egress Docker lab (\texttt{gitea/gitea:1.27.0} against \texttt{1.27.1},
one variable at a time) with a canary-first payload that writes one file and
does nothing else, we contribute:

\begin{enumerate}
\item \textbf{Process-level proof of execution}, which the advisory does not
provide. Our hook reports its own lineage: parent process
\texttt{/usr/bin/git write-tree}, own command line
\texttt{/bin/sh\hspace{0pt}hooks/post-index-change\hspace{0pt}0\hspace{0pt}0}, working
directory
\texttt{/data/gitea/tmp/\allowbreak local-repo/\allowbreak upload.git749828539},
\texttt{GIT\_DIR=.}, \texttt{uid=1000(git)}, plus the container's process tree
captured from inside the hook (Section~\ref{sec:repro}). The public record says
generically that Git invokes the hook ``while writing the index''; in the public
sources we checked (advisory, CVE record, fix commit message, patch diff) the
parent process is not named.
\item \textbf{Empirical verification of preconditions the vendor only asserts},
and a derived finding that changes triage: in the official image as published,
all three hold \emph{without setting any precondition-related option}
(Section~\ref{sec:precond}).
\item \textbf{Fix verification with a live negative}, not an absence of
evidence. On \texttt{1.27.1} with byte-identical payload: no execution in 22 of
22 cycles, the second \texttt{diffpatch} submission changed from HTTP 201 to
HTTP 500, and a polling monitor still observed the payload materialising at the
sink path --- so the arm proves the payload arrives and still does not run
(Section~\ref{sec:fix}).
\item \textbf{The non-obvious consequence of the fix}, which is the most useful
thing here for a defender. The decisive hunk changes
\texttt{Clone(\dots,\allowbreak true)} to \texttt{Clone(\dots,\allowbreak false)}
and adds no path deny-list, no mode filter, and no change to the
\texttt{git apply} flags. The same patch bytes are still accepted and still land
on disk after the fix. Consequence: patch-content detection cannot distinguish
an attack that executed from one that was neutralised, and an artifact rule
fires on both (Sections~\ref{sec:fixchange} and~\ref{sec:detect}).
\item \textbf{Measured impact radius} in place of an enumerated list:
\texttt{app.ini} readable (1323 bytes, 3 secret-bearing lines), planted secrets
read back, arbitrary write confirmed where the service lives, repositories
visible --- and one honest negative, the parent process environment leaked
nothing (Section~\ref{sec:impact}).
\item \textbf{A negative-results log}. Four of our own measurements were wrong
or unprovable and they are printed in Section~\ref{sec:negative}, including a
\texttt{grep} that appeared to prove the \texttt{diffpatch} route was absent
because it inspected a 581-byte wrapper instead of the 124\,MB binary. A
published exploit does not make the surrounding evidence trustworthy by
itself.
\end{enumerate}

Section~\ref{sec:mech} gives the mechanism plus an independent static
corroboration; Section~\ref{sec:repro} the reproduction and process evidence;
Section~\ref{sec:fix} the A/B; Section~\ref{sec:detect} what a defender can
actually see; Section~\ref{sec:remediation} remediation and incident response;
Section~\ref{sec:limitations} limitations and negative results; and
Section~\ref{sec:ethics} ethics, scope and reproducibility. Every number was
either measured in our lab or read from the public record with a date. Nothing
here is estimated.

% =============================================================================
\section{The Public Record, Ethics Position, and Timeline}
\label{sec:record}

\subsection{What the vendor published}

Table~\ref{tab:record} separates the two bodies of work. Nothing in this paper
moves from the right column to the left.

\begin{table*}[t]
\centering
\footnotesize
\begin{tabular}{@{}p{7.7cm}p{8.1cm}@{}}
\toprule
\textbf{Public record (vendor and registries)} & \textbf{This work} \\
\midrule
Mechanism: add/add collision, three-way fallback, bare root as
\texttt{\$GIT\_DIR}, live hook (GHSA ``Details''). &
Process-level proof of that mechanism: parent PID and command line, hook
command line, \texttt{cwd}, \texttt{GIT\_DIR}, resolved hook path, container
process tree, wall-clock timing. \\
Vulnerable command \texttt{git apply --index --recount --cached --binary} plus
\texttt{-3} on Git $\geq$2.32, in
\texttt{services/repository/files/patch.go}. &
Independent static corroboration of the collision detector recovered from the
shipped binary (\texttt{git ls-files -u -z}, Section~\ref{sec:mech}), and the
measurement that Git is 2.54.0 in both images. \\
Working exploit \texttt{gitea\_diffpatch\_rce\_poc.py} ($\approx$330 lines) with
output-free exfiltration via Git objects plus a branch fetched over
authenticated smart HTTP. &
A deliberately weaker stdlib-only canary PoC: no outbound channel, no reverse
shell, no destructive action; the hook writes one canary file with process
evidence (Section~\ref{sec:repro}). We did \emph{not} reproduce the vendor's
exfiltration channel. \\
Three stated preconditions: Git $\geq$2.32, enabled \texttt{diffpatch} route,
writable and executable temporary filesystem. &
Those three measured rather than assumed, plus the derived finding that the
official image satisfies all three by default
(Section~\ref{sec:precond}). \\
Impact stated as a list: application secrets, database credentials, mounted
repositories, OAuth and integration credentials. &
Impact measured: byte counts, secret-bearing line counts, planted secrets read
back, arbitrary-write confirmation, network visibility, and one negative
(Section~\ref{sec:impact}). \\
Fix commit \texttt{d7bc52be\ldots} (2026-07-26) and first patched version
1.27.1. &
A/B execution of the fix with a live negative, plus a response-code change
(HTTP 500 on the second submission) that the advisory does not mention
(Section~\ref{sec:fix}). \\
\bottomrule
\end{tabular}
\caption{Credit division. Everything in the left column belongs to the vendor
and the registries; we cite it and do not claim it. The right column is what
this paper adds.}
\label{tab:record}
\end{table*}

\subsection{Timeline and operational urgency}
\label{sec:timeline}

\begin{figure*}[t]
\centering
\includegraphics[width=\textwidth]{\figdir/fig-timeline.png}
\caption{Disclosure and response chronology, read from the registries rather
than remembered: fix commit merged 2026-07-26T17:32:02Z (three days
\emph{before} the advisory), GHSA published 2026-07-28T17:34:46Z and updated
18:40:54Z the same day, CISA KEV \texttt{dateAdded} 2026-08-25 with FCEB
\texttt{dueDate} 2026-08-28, and NVD publication 2026-08-26T20:17:56.010Z
with status \emph{Analyzed}. The KEV note references BOD 26-04, whose
staggered windows are 3, 14 and 60 days.}
\label{fig:timeline}
\end{figure*}

Two facts in Figure~\ref{fig:timeline} make this urgent rather than academic.
First, the fix reached the repository on 2026-07-26, three days before the
advisory: an operator who tracks commits had a patch before a disclosure
notice. Second, the response clock arrived a month later and is short --- CISA
added the CVE to the KEV catalog on 2026-08-25, 28 days after the advisory,
NVD published it the following day with \emph{Analyzed} status, and the federal
(FCEB) deadline is 2026-08-28, three days after the KEV entry. In this
case the operational urgency preceded the NVD record by a day. The KEV entry's own notes reference BOD 26-04, with staggered
windows of 3, 14 and 60 days.

Two catalog measurements put that deadline in proportion. The KEV record
shipped with this paper was fetched on 2026-08-27 (catalog
\texttt{2026.08.27}, 1685 entries in 1\,618\,253 bytes, full-file SHA-256
pinned in the supporting material): 1672 entries (99.2\,\%) already had a
passed \texttt{dueDate} and only 13 were due within 30 days. A KEV entry is a
prioritised signal, not a self-enforcing control. The same entry records
\texttt{knownRansomwareCampaignUse = Unknown}, so we do not claim ransomware
use, and NVD records no CVSS v4.0 score, so we do not quote one. The severity
of record is the NVD-measured CVSS 3.1 vector
\texttt{AV:N/\hspace{0pt}AC:L/\hspace{0pt}PR:N/\hspace{0pt}UI:N/\hspace{0pt}S:U/\hspace{0pt}C:H/\hspace{0pt}I:H/\hspace{0pt}A:H}
= 9.8 (Critical), typed \emph{Secondary}, with CWE-94.

\subsection{Scope statement}

Every experiment ran in a Docker lab we created and destroyed: an
\texttt{--internal} bridge with verified zero egress, a fictional hostname, no
real users and no real data. Payloads were canary-first by policy.
Section~\ref{sec:ethics} gives the full ethics, consent and reproducibility
statement rather than repeating it here.
